\begin{example}[不能用单个幺正矩阵对角化的 Yukawa 耦合]
取两代模型
\begin{equation*}
Y=\frac1{\sqrt2}\begin{pmatrix}2&-\ii\\-2&-\ii\end{pmatrix},\qquad
U_L=\frac1{\sqrt2}\begin{pmatrix}1&1\\-1&1\end{pmatrix},\qquad
U_R=\operatorname{diag}(1,\ii).
\end{equation*}
求质量本征值并检查 Higgs 顶角是否保留味非对角项。
\end{example}
\begin{proof}[解]
直接相乘得到
\begin{align*}
YY^\dagger&=\begin{pmatrix}5/2&-3/2\\-3/2&5/2\end{pmatrix},
&Y^\dagger Y&=\begin{pmatrix}4&0\\0&1\end{pmatrix},\\
U_L^\dagger YU_R&=\operatorname{diag}(2,1).
\end{align*}
左右两个厄米矩阵不相等，所以 $Y$ 不是正规矩阵，无法作幺正相似对角化。
奇异值分解却给实正质量 $m_1c^2=\sqrt2v_E$、$m_2c^2=v_E/\sqrt2$。
由于质量项与 $h$ 顶角都乘同一个 $Y$，相同的左右旋转给
\begin{equation*}
\Lag_{hff}=-\frac{\sqrt{\hbar c}\,h}{v_E}
\left(m_1c^2\bar f_1'f_1'+m_2c^2\bar f_2'f_2'\right).
\end{equation*}
规范基中的复非对角系数并没有产生树级 Higgs 味改变。
若引入第二个独立标量的 Yukawa 矩阵，这个共同对角化条件一般就不成立。
\end{proof}
